% =====================================================================
\chapter{Variables aleatorias continuas y distribución normal}
\label{tema:3}
% =====================================================================

Las magnitudes que varían en un continuo, como los pesos de una red o el número aleatorio que se usa para muestrear, se describen con densidades. Este tema introduce las variables continuas, estudia la uniforme, la exponencial y la normal, demuestra el método de la transformada inversa y aplica las propiedades de la varianza a la inicialización de redes neuronales. Se sigue el capítulo~3 y la §4.1 de~\cite{bt}, la unidad~II de~\cite{mit6041}, las clases~5b a~6a de~\cite{mit1805} y el capítulo~5 de~\cite{bh}; la normal en aprendizaje automático, en~\cite[§6.5]{mml} y~\cite[§1.2.4]{bishop}, y la inicialización, en los artículos originales~\cite{glorot,he}.

\section{Densidad y función de distribución}

\begin{definicion}[Variable continua]\label{def:continua}
Una variable aleatoria $X$ es \emph{continua} si existe una función $f_X\ge 0$, su \emph{densidad}, tal que
\[
\PP(a\le X\le b)=\int_a^b f_X(x)\,dx\qquad\text{para todo } a\le b,
\]
y en particular $\int_{-\infty}^{\infty}f_X(x)\,dx=1$. La \emph{función de distribución} de cualquier variable es $F_X(x)=\PP(X\le x)$.
\end{definicion}

\begin{proposicion}\label{prop:densidad}
Si $X$ es continua con densidad $f_X$:
\begin{enumerate}
  \item $\PP(X=x)=0$ para todo $x$; por tanto, da igual usar desigualdades estrictas o no estrictas;
  \item $F_X(x)=\int_{-\infty}^{x}f_X(t)\,dt$ y, en cada punto de continuidad de $f_X$, $F_X'(x)=f_X(x)$;
  \item $F_X$ es no decreciente, con $\lim_{x\to-\infty}F_X(x)=0$ y $\lim_{x\to\infty}F_X(x)=1$.
\end{enumerate}
\end{proposicion}

\begin{proof}
(i) $\PP(X=x)=\int_x^x f_X=0$. (ii) La primera igualdad es la definición con $a\to-\infty$; la segunda es el teorema fundamental del cálculo. (iii) Si $x<y$, $F_X(y)-F_X(x)=\int_x^y f_X\ge 0$; los límites se siguen de que la integral total es~1.
\end{proof}

La densidad no es una probabilidad: puede ser mayor que~1. Para $h$ pequeño, $\PP(x\le X\le x+h)\approx f_X(x)\,h$, de modo que $f_X(x)$ es probabilidad por unidad de longitud. La esperanza y la varianza se definen como en el caso discreto cambiando sumas por integrales: $\E[g(X)]=\int g(x)f_X(x)\,dx$, y todas las propiedades del \cref{tema:2} se mantienen~\cite[§3.1]{bt}.

\section{Uniforme y exponencial}

\begin{definicion}[Uniforme]
$X\sim\Unif(a,b)$, con $a<b$, si $f_X(x)=1/(b-a)$ en $[a,b]$ y $0$ fuera.
\end{definicion}

\begin{proposicion}\label{prop:uniforme}
Si $X\sim\Unif(a,b)$, entonces $\E[X]=(a+b)/2$ y $\Var(X)=(b-a)^{2}/12$.
\end{proposicion}

\begin{proof}
$\E[X]=\frac{1}{b-a}\int_a^b x\,dx=\frac{b^{2}-a^{2}}{2(b-a)}=\frac{a+b}{2}$. Para la varianza, la variable $Y=X-\frac{a+b}{2}$ es uniforme en $[-c,c]$ con $c=\frac{b-a}{2}$, y $\Var(X)=\Var(Y)=\E[Y^{2}]=\frac{1}{2c}\int_{-c}^{c}y^{2}\,dy=\frac{c^{2}}{3}=\frac{(b-a)^{2}}{12}$.
\end{proof}

\begin{definicion}[Exponencial]
$X\sim\Expo(\lambda)$, con $\lambda>0$, si $f_X(x)=\lambda e^{-\lambda x}$ para $x\ge 0$ y $0$ para $x<0$. Su función de distribución es $F_X(x)=1-e^{-\lambda x}$ para $x\ge 0$.
\end{definicion}

Integrando por partes se obtiene $\E[X]=1/\lambda$ y $\E[X^{2}]=2/\lambda^{2}$, luego $\Var(X)=1/\lambda^{2}$~\cite[§3.1]{bt}. La exponencial modela tiempos de espera sin desgaste, como el de un taxi o el de una desintegración radiactiva~\cite[clase~5c]{mit1805}.

\begin{proposicion}[Falta de memoria]\label{prop:memoria}
Si $X\sim\Expo(\lambda)$ y $s,t\ge 0$, entonces $\PP(X>s+t\mid X>s)=\PP(X>t)$.
\end{proposicion}

\begin{proof}
Como $\{X>s+t\}\subseteq\{X>s\}$, $\PP(X>s+t\mid X>s)=\dfrac{\PP(X>s+t)}{\PP(X>s)}=\dfrac{e^{-\lambda(s+t)}}{e^{-\lambda s}}=e^{-\lambda t}=\PP(X>t)$.
\end{proof}

\section{La distribución normal}

\begin{definicion}[Normal]\label{def:normal}
$X\sim\N(\mu,\sigma^{2})$, con $\mu\in\R$ y $\sigma>0$, si
\[
f_X(x)=\frac{1}{\sqrt{2\pi\sigma^{2}}}\exp\Bigl(-\frac{(x-\mu)^{2}}{2\sigma^{2}}\Bigr).
\]
La \emph{normal estándar} es $\N(0,1)$; su función de distribución se denota $\Phi$, y el valor máximo de su densidad es $1/\sqrt{2\pi}\approx 0{,}399$.
\end{definicion}

Que la densidad integre~1 se deduce de la integral gaussiana $\int_{-\infty}^{\infty}e^{-x^{2}/2}\,dx=\sqrt{2\pi}$, que se calcula pasando a coordenadas polares en el plano~\cite[§3.3]{bt}. Por simetría, $\E[X]=\mu$, y una integración por partes da $\Var(X)=\sigma^{2}$.

\begin{proposicion}[Transformaciones lineales]\label{prop:afin}
Si $X\sim\N(\mu,\sigma^{2})$ y $a\ne 0$, $b$ son constantes, entonces $aX+b\sim\N(a\mu+b,\,a^{2}\sigma^{2})$. En particular, $Z=(X-\mu)/\sigma\sim\N(0,1)$ y $\PP(X\le x)=\Phi\bigl((x-\mu)/\sigma\bigr)$.
\end{proposicion}

\begin{proof}
Sea $a>0$ (el caso $a<0$ es análogo). Para $Y=aX+b$, $F_Y(y)=\PP(X\le (y-b)/a)=F_X\bigl((y-b)/a\bigr)$. Derivando con la regla de la cadena,
\[
f_Y(y)=\frac1a f_X\Bigl(\frac{y-b}{a}\Bigr)=\frac{1}{\sqrt{2\pi a^{2}\sigma^{2}}}\exp\Bigl(-\frac{(y-b-a\mu)^{2}}{2a^{2}\sigma^{2}}\Bigr),
\]
que es la densidad de $\N(a\mu+b,a^{2}\sigma^{2})$.
\end{proof}

Gracias a la \cref{prop:afin}, cualquier probabilidad normal se reduce a valores de $\Phi$, que se tabulan o se calculan numéricamente. El \cref{tab:normal} recoge los más usados y la \cref{fig:normal} los representa.

\begin{table}[h]
\centering\small
\caption{Probabilidad de que una normal caiga a menos de $k$ desviaciones típicas de su media.}\label{tab:normal}
\begin{tabular}{@{}lcccc@{}}
\toprule
$k$ & $1$ & $1{,}96$ & $2$ & $3$ \\
\midrule
$\PP(\abs{X-\mu}\le k\sigma)=2\Phi(k)-1$ & $0{,}683$ & $0{,}950$ & $0{,}954$ & $0{,}997$ \\
\bottomrule
\end{tabular}
\end{table}

\begin{figure}[h]
\centering
\begin{tikzpicture}[x=1.15cm,y=6cm,font=\small]
  \fill[gray!35,domain=-1:1,samples=40] (-1,0) -- plot (\x,{exp(-\x*\x/2)/sqrt(2*pi)}) -- (1,0) -- cycle;
  \fill[gray!15,domain=1:2,samples=20] (1,0) -- plot (\x,{exp(-\x*\x/2)/sqrt(2*pi)}) -- (2,0) -- cycle;
  \fill[gray!15,domain=-2:-1,samples=20] (-2,0) -- plot (\x,{exp(-\x*\x/2)/sqrt(2*pi)}) -- (-1,0) -- cycle;
  \draw[-{Stealth[length=2mm]}] (-3.6,0) -- (3.7,0) node[right] {$z$};
  \draw[thick,domain=-3.5:3.5,samples=80] plot (\x,{exp(-\x*\x/2)/sqrt(2*pi)});
  \foreach \k in {-3,-2,-1,0,1,2,3} {\draw (\k,0.012) -- (\k,-0.012) node[below] {$\k$};}
  \node at (0,0.17) {$0{,}683$};
  \draw[{Stealth[length=1.5mm]}-{Stealth[length=1.5mm]}] (-2,0.47) -- node[above] {$0{,}954$} (2,0.47);
  \draw[dotted] (-2,0) -- (-2,0.47) (2,0) -- (2,0.47);
\end{tikzpicture}
\caption{Densidad de la normal estándar. La zona oscura, $\abs{z}\le 1$, contiene el $68{,}3\,\%$ de la probabilidad; junto con las zonas claras, $\abs{z}\le 2$, el $95{,}4\,\%$.}
\label{fig:normal}
\end{figure}

\begin{ejemplo}\label{ej:normal}
Si $Z\sim\N(0,1)$, $\PP(-1{,}5\le Z\le 2)=\Phi(2)-\Phi(-1{,}5)=0{,}9772-0{,}0668=0{,}9104$~\cite[clase~5c]{mit1805}.
\end{ejemplo}

La normal aparece cuando una magnitud es la suma de muchos efectos pequeños e independientes; el teorema central del límite (\cref{tema:5}) da la justificación precisa.

\section{Muestreo por transformada inversa}

Un ordenador genera números (pseudo)aleatorios uniformes en $(0,1)$. El resultado siguiente permite convertirlos en muestras de cualquier distribución continua cuya función de distribución sea invertible; es la versión continua de la \cref{prop:inversa-discreta}.

\begin{teorema}[Universalidad de la uniforme]\label{teo:inversa}
Sea $F$ una función de distribución continua y estrictamente creciente en el intervalo donde $0<F<1$, y sea $U\sim\Unif(0,1)$. Entonces $X=F^{-1}(U)$ tiene función de distribución $F$.
\end{teorema}

\begin{proof}
Como $F$ es estrictamente creciente, $F^{-1}(u)\le x$ si y solo si $u\le F(x)$. Por tanto, $\PP(X\le x)=\PP(U\le F(x))=F(x)$, porque $\PP(U\le c)=c$ para $c\in[0,1]$.
\end{proof}

\begin{observacion}
El \cref{teo:inversa} es la base de la generación de muestras en todas las bibliotecas numéricas. Para la normal, $\Phi^{-1}$ no tiene expresión cerrada y se usan otros métodos (por ejemplo, el de Box--Muller), pero el principio es el mismo: transformar uniformes~\cite[§6.7]{mml}. El \cref{eje:3-exponencial} lo aplica a la exponencial.
\end{observacion}

\section{Aplicación: inicialización de redes neuronales}

Una neurona lineal calcula $y=\sum_{j=1}^{n}w_jx_j$, donde $n$ es el número de entradas (el \emph{fan-in}). Antes de entrenar, los pesos $w_j$ se eligen al azar. ¿Con qué varianza?

\begin{proposicion}[Varianza de una neurona]\label{prop:neurona}
Supóngase que $w_1,\dots,w_n,x_1,\dots,x_n$ son independientes, que los $w_j$ tienen media~0 y varianza $\Var(w)$ y que los $x_j$ tienen media~0 y varianza $\Var(x)$. Entonces
\[
\E[y]=0,\qquad \Var(y)=n\,\Var(w)\,\Var(x).
\]
\end{proposicion}

\begin{proof}
Por independencia, $\E[w_jx_j]=\E[w_j]\E[x_j]=0$, luego $\E[y]=0$. Además, $\Var(w_jx_j)=\E[w_j^{2}x_j^{2}]-0=\E[w_j^{2}]\,\E[x_j^{2}]=\Var(w)\Var(x)$, porque ambas medias son nulas. Los $n$ productos $w_jx_j$ son independientes (cada uno depende de variables distintas), y la \cref{prop:suma} da $\Var(y)=\sum_j\Var(w_jx_j)=n\Var(w)\Var(x)$.
\end{proof}

Si se toma $\Var(w)=1$ y $n=512$, cada capa multiplica la varianza por~512 y la desviación típica por $\sqrt{512}\approx 22{,}6$; tras diez capas la desviación ha crecido en un factor $512^{5}\approx 3{,}5\cdot 10^{13}$ y la red se desborda. Si $\Var(w)$ es demasiado pequeña, la señal se anula. La \cref{prop:neurona} da la elección que conserva la varianza:
\[
\Var(w)=\frac1n,\qquad\text{por ejemplo } n=512\ \Rightarrow\ \sigma_w=1/\sqrt{512}\approx 0{,}0442.
\]
Glorot y Bengio~\cite{glorot} propusieron este análisis y, para equilibrar también la propagación hacia atrás de los gradientes, la elección de compromiso $\Var(w)=2/(n_{\mathrm{ent}}+n_{\mathrm{sal}})$. Con la activación ReLU, $\max(0,x)$, la mitad de las entradas se anula y el segundo momento se reduce a la mitad (\cref{eje:3-he}); He \emph{et al.}~\cite{he} compensan con $\Var(w)=2/n$, que para $n=512$ da $\sigma_w=\sqrt{2/512}=0{,}0625$. El módulo~9 retoma estas reglas.

\begin{lecturas}
\item Bertsekas y Tsitsiklis~\cite{bt}, §3.1 «Continuous Random Variables and PDFs», §3.2 «Cumulative Distribution Functions», §3.3 «Normal Random Variables» y §4.1 «Derived Distributions».
\item MIT 6.041SC~\cite{mit6041}, unidad~II, lecciones~8 «Continuous Random Variables» y~10 «Continuous Bayes' Rule; Derived Distributions».
\item MIT 18.05~\cite{mit1805}, clases~5b «Continuous Random Variables», 5c «Gallery of Continuous Random Variables» y 6a (esperanza y varianza en el caso continuo).
\item Blitzstein y Hwang~\cite{bh}, cap.~5 «Continuous Random Variables» (universalidad de la uniforme).
\item Deisenroth, Faisal y Ong~\cite{mml}, §6.5 «Gaussian Distribution» y §6.7 «Change of Variables/Inverse Transform».
\item Glorot y Bengio~\cite{glorot}, y He \emph{et al.}~\cite{he}, sección dedicada a la inicialización.
\end{lecturas}

\section*{Ejercicios del tema 3}
\addcontentsline{toc}{section}{Ejercicios del tema 3}

\begin{ejercicio}{1}{eje:3-tabla}
Sea $X\sim\N(10,4)$. Usando solo el \cref{tab:normal} y la simetría de la normal, calcule $\PP(8\le X\le 12)$ y $\PP(X>14)$.
\end{ejercicio}

\begin{ejercicio}{2}{eje:3-espera}
El tiempo de espera $X$, en minutos, hasta la siguiente petición a un servidor sigue una $\Expo(0{,}5)$. (a) Calcule $\E[X]$ y $\sigma_X$. (b) Calcule $\PP(X>2)$. (c) Si ya se ha esperado 3 minutos, ¿qué probabilidad hay de esperar más de 5 en total?
\end{ejercicio}

\begin{ejercicio}{2}{eje:3-cuadrado}
Sea $U\sim\Unif(-1,1)$. (a) Calcule $\E[U]$, $\Var(U)$ y $\E[U^{4}]$. (b) Halle la función de distribución y la densidad de $U^{2}$. ¿Es $U^{2}$ uniforme en $(0,1)$?
\fuente{adaptado de Harvard Stat 110~\cite{stat110}, \emph{Homework}~5, problema~3.}
\end{ejercicio}

\begin{ejercicio}{2}{eje:3-init}
Una capa lineal tiene $n=4096$ entradas de media~0 y varianza~1. (a) ¿Qué desviación típica de los pesos conserva la varianza (regla de fan-in)? (b) ¿Y con la regla de He? (c) Si se usa $\sigma_w=0{,}1$ con $n=256$, ¿cuánto vale $\Var(y)$ y qué ocurre tras diez capas?
\end{ejercicio}

\begin{ejercicio}{3}{eje:3-exponencial}
Sea $U\sim\Unif(0,1)$ y $\lambda>0$. Construya, como función de $U$, una variable con distribución $\Expo(\lambda)$ y demuestre que la tiene. Muestre además que $-\ln(U)/\lambda$ también es una $\Expo(\lambda)$.
\fuente{adaptado de Harvard Stat 110~\cite{stat110}, \emph{Strategic Practice}~5, sección~3, problema~2.}
\end{ejercicio}

\begin{ejercicio}{3}{eje:3-he}
Sea $x$ una variable continua con densidad simétrica respecto de~0 y varianza $v$, y sea $r=\max(0,x)$. (a) Demuestre que $\E[r^{2}]=v/2$. (b) En una neurona $y=\sum_{j=1}^{n}w_jr_j$, con $w_j$ de media~0 independientes de los $r_j$, demuestre que $\Var(y)=n\Var(w)\,\E[r^{2}]$ aunque $\E[r_j]\ne 0$. (c) Deduzca que $\Var(w)=2/n$ hace $\Var(y)=v$.
\end{ejercicio}
